El problema se centra en el an\'alisis del conjunto de datos Boston, disponible en \texttt{MASS}, con el objetivo de describir su estructura, identificar patrones relevantes y responder las preguntas de investigaci\'on formuladas en la tarea. La metodolog\'ia se organiza seg\'un CRISP-DM en las fases de comprensi\'on del negocio, comprensi\'on de los datos, preparaci\'on, modelado y evaluaci\'on, dejando la implementaci\'on computacional como soporte reproducible del estudio.

El objetivo general es caracterizar los factores asociados al valor mediano de las viviendas y explorar si el PCA ayuda a sintetizar la informaci\'on multivariada. Como objetivos espec\'ificos se busca detectar at\'ipicos, identificar los suburbios con menor valor de vivienda, cuantificar la relaci\'on entre el n\'umero de habitaciones y el precio, evaluar el efecto de \textit{chas} y \textit{lstat}, y revisar si la tasa de criminalidad puede explicarse a partir de otras variables. El documento se organiza en fundamentos te\'oricos, desarrollo metodol\'ogico, discusi\'on, conclusiones y anexo.